\chapter{Resultados}
\label{chap:resultados}

\section{QA1 — Automatização do Protocolo}
\label{sec:res_qa1}

O protocolo record--replay foi executado de forma automática em duas etapas.
Primeiro, uma baseline limpa foi gravada na rota
\texttt{dissertation\_\allowbreak clean01}, gerando o arquivo YAML de rota e um
bag MCAP com \texttt{/\allowbreak scan}, \texttt{/\allowbreak odom},
\texttt{/\allowbreak imu} e \texttt{/\allowbreak pose}. Em seguida, foram
executados dez replays independentes dessa rota. Para cada replay, um script de
campanha relançou completamente a simulação, a pilha Nav2, o orquestrador e o
coletor de dados, aguardou a ativação dos nós lifecycle do Nav2 e então iniciou
a coleta e a reprodução da rota.

A campanha final foi concluída com dez replays válidos, todos com
\texttt{success: true}. Os artefatos principais foram armazenados em
\texttt{fleet\_\allowbreak ws/\allowbreak runs/\allowbreak dissertation\_\allowbreak clean01\_\allowbreak final\_\allowbreak manual/},
incluindo o manifesto da campanha, os exports individuais de cada replay, os
bags MCAP, os CSVs de trajetória, o gráfico de sobreposição das trajetórias e o
arquivo \texttt{summary.json} consolidado.

\section{QA2 — RMSE de Repetibilidade}
\label{sec:res_qa2}

A Tabela~\ref{tab:campanha_final} apresenta as métricas por replay da campanha
final. O primeiro replay é usado como execução de referência; por isso, seu RMSE
e erro final contra a referência são iguais a zero e não entram no cálculo das
médias de erro.

\begin{table}[H]
\centering
\caption{Campanha final controlada: duração, comprimento, RMSE e erro final por replay.}
\label{tab:campanha_final}
\footnotesize
\begin{tabular}{lrrrr}
\toprule
\textbf{Execução} & \textbf{Duração (s)} & \textbf{Comprimento (m)} & \textbf{RMSE vs.\ ref. (cm)} & \textbf{Erro final (cm)} \\
\midrule
Replay 01 & 17,96 & 0,840 & 0,00 & 0,00 \\
Replay 02 & 17,64 & 0,807 & 2,57 & 3,26 \\
Replay 03 & 18,11 & 0,823 & 2,01 & 1,64 \\
Replay 04 & 17,57 & 0,824 & 2,58 & 1,56 \\
Replay 05 & 17,42 & 0,825 & 3,77 & 1,47 \\
Replay 06 & 17,32 & 0,847 & 4,58 & 0,67 \\
Replay 07 & 17,28 & 0,826 & 4,19 & 1,36 \\
Replay 08 & 17,32 & 0,825 & 4,76 & 1,50 \\
Replay 09 & 17,71 & 0,830 & 1,65 & 1,00 \\
Replay 10 & 16,78 & 0,817 & 4,08 & 2,26 \\
\bottomrule
\end{tabular}
\end{table}

Considerando os nove replays comparados à referência, o RMSE médio foi de
3,35\,cm, com desvio-padrão de 1,10\,cm. O menor RMSE observado foi de
1,65\,cm e o maior foi de 4,76\,cm. Todos os valores ficaram substancialmente
abaixo da tolerância padrão de objetivo do Nav2 de 25\,cm.

O resultado é metodologicamente mais forte que as campanhas preliminares porque
os pontos iniciais foram controlados. No \texttt{summary.json}, todos os dez
replays apresentam \texttt{start\_\allowbreak xy} praticamente igual a $(0,0)$,
com desvio máximo da ordem de $10^{-18}$\,m em relação ao primeiro replay.
Portanto, o RMSE medido não está misturando diferenças progressivas de pose
inicial com erro de repetição da trajetória.

A distância final em relação ao \texttt{replay\_\allowbreak 01} também permaneceu
baixa: o erro final médio foi de 1,64\,cm, com máximo de 3,26\,cm. Isso indica
que, além de seguir trajetórias semelhantes ao longo do percurso, o robô terminou
consistentemente em posições próximas sob o protocolo controlado.

\section{QA3 — Consistência Temporal}
\label{sec:res_qa3}

A duração média dos dez replays foi de 17,51\,s, com desvio-padrão de
0,36\,s. O comprimento médio estimado por \texttt{/\allowbreak odom} foi de
0,826\,m, com desvio-padrão de 0,010\,m. A variação temporal observada é
pequena para o objetivo da análise e compatível com o uso de reamostragem em
tempo normalizado.

Esses valores mostram que as execuções foram não apenas espacialmente próximas,
mas também temporalmente consistentes. A combinação de baixo RMSE, baixo erro
final e baixa dispersão de duração sustenta a interpretação de que o framework
reproduziu a rota de forma repetível sob condições controladas.

\section{Integridade dos Bags de Sensores}
\label{sec:res_bags}

Os dez bags da campanha final contêm os tópicos planejados:
\texttt{/\allowbreak scan}, \texttt{/\allowbreak odom}, \texttt{/\allowbreak imu}
e \texttt{/\allowbreak pose}. A Tabela~\ref{tab:sensor_bags} resume as
contagens médias de mensagens e as frequências efetivas estimadas por tópico.

\begin{table}[H]
\centering
\caption{Frequências efetivas médias de coleta de sensores na campanha final.}
\label{tab:sensor_bags}
\begin{tabular}{lrrr}
\toprule
\textbf{Tópico} & \textbf{Msgs/exec. média} & \textbf{Faixa de msgs} & \textbf{Freq. efetiva média (Hz)} \\
\midrule
\texttt{/\allowbreak scan} & 175,4 & 168--181 & 10,0 \\
\texttt{/\allowbreak odom} & 487,4 & 467--504 & 27,8 \\
\texttt{/\allowbreak imu}  & 3504,7 & 3364--3630 & 199,7 \\
\texttt{/\allowbreak pose} & 4,9 & 4--7 & 0,3 \\
\bottomrule
\end{tabular}
\end{table}

A trajetória principal da análise foi extraída de \texttt{/\allowbreak odom};
\texttt{/\allowbreak pose} foi preservado como dado auxiliar. A presença
consistente desses tópicos nos dez bags confirma que a coleta em MCAP permaneceu
operacional mesmo com o relançamento completo da simulação entre replays.

\section{Experimentos Complementares: Variabilidade Fora do Protocolo Controlado}
\label{sec:res_exp_complementares}

Antes da campanha final controlada, foram executadas campanhas exploratórias nas
quais os replays não eram reinicializados a partir do mesmo estado. Esses dados
foram úteis para diagnosticar uma ameaça metodológica: o início de uma execução
podia coincidir com o fim da anterior, especialmente quando
\texttt{return\_\allowbreak to\_\allowbreak start=false}. Nesses casos, a abertura
entre curvas não pode ser interpretada diretamente como erro puro de
repetibilidade, pois mistura repetição da trajetória com variação de pose inicial.

Um exemplo desse problema apareceu em campanhas nas quais o \texttt{start\_\allowbreak xy}
de um replay era praticamente igual ao \texttt{end\_\allowbreak xy} do replay
anterior. O efeito visual era uma separação progressiva entre trajetórias, mas
essa separação não media apenas a variabilidade do controlador ou do SLAM; media
também a consequência de iniciar cada nova execução a partir de uma pose diferente.

Por essa razão, os resultados exploratórios são mantidos como evidência
diagnóstica e motivação do protocolo final, mas não substituem a campanha
\texttt{dissertation\_\allowbreak clean01\_\allowbreak final\_\allowbreak manual}, na qual
o estado inicial foi efetivamente controlado. O principal valor desses dados é
mostrar empiricamente que a infraestrutura de coleta não basta: para que o RMSE
seja interpretável como repetibilidade de trajetória, o protocolo precisa impor
mesma rota, mesma ação de navegação, mesma condição inicial e coleta padronizada.

\section{Considerações Finais}
\label{sec:res_conclusao}

Os resultados respondem afirmativamente às três primeiras questões de avaliação.
QA1 é respondida pela execução automatizada da baseline e dos dez replays, com
geração automática de bags, CSVs, gráfico e \texttt{summary.json}. QA2 é
respondida pelo RMSE médio de 3,35\,cm e máximo de 4,76\,cm contra a execução de
referência, ambos muito abaixo da tolerância de 25\,cm do Nav2. QA3 é respondida
pela baixa variabilidade temporal observada nos dez replays, com duração média
de 17,51\,s e desvio-padrão de 0,36\,s.

A análise também reforça uma lição metodológica central: a pose inicial precisa
ser controlada explicitamente. Campanhas preliminares sem reinicialização entre
replays produziram métricas contaminadas por deslocamentos iniciais progressivos.
A campanha final evita essa ameaça ao relançar a simulação e a pilha de navegação
antes de cada replay, fazendo com que as dez trajetórias analisadas partam do
mesmo estado em \texttt{/\allowbreak odom}. O capítulo seguinte conclui a
dissertação, discutindo implicações, limitações e trabalhos futuros.
